\documentclass[11pt]{letter}
\usepackage[margin=1in]{geometry}
\usepackage{hyperref}
\usepackage{amsmath}

\signature{Tan XIN\\Independent Researcher, IAIP\\Email: summicron352@gmail.com}

\address{Independent Researcher\\International Association of Independent Physicists (IAIP)}

\date{\today}

\begin{document}

\begin{letter}{Editorial Team\\The Astrophysical Journal}

\opening{Dear Editorial Team of The Astrophysical Journal,}

I am submitting the manuscript entitled ``Surface Density as the Organizing Principle of Little Red Dot Spectral Properties'' for consideration as a full research article in \emph{The Astrophysical Journal} (ApJ).

\textbf{Scientific Significance:}

Currently, JWST's Little Red Dots (LRDs) pose a major puzzle, as their extreme compactness, anomalous reddening, broad emission lines, and mid-infrared silence challenge standard galaxy evolution models. This manuscript provides a unified solution aimed at resolving the current dilemmas facing these paradigms. By leveraging a catalog of 296 credible, public LRD sources to screen and exhaustively analyze 1,626 matched pairs, I rigorously demonstrate that surface density ($\Sigma$) is the fundamental, independent organizing parameter of these anomalies (supported by a $5.6\sigma$ partial correlation and a 34.6\% color-reversal rate, $p < 10^{-30}$). These findings are independently corroborated by Euclid compact-source candidates and JADES DR5 deep-field cross-validation. Building on this robust empirical foundation, I propose a mechanism-agnostic effective gravity framework, $G_{\rm eff}(\Sigma)$, that quantitatively unifies all four hallmark LRD features through deepened potential wells.

\textbf{Manuscript History:}

A shorter version of this work was previously submitted to ApJL (Manuscript ID: AAS76108). The ApJL Editor found the topic interesting but concluded that the mechanism-focused nature of the draft was better suited for a full-length format rather than a rapid-communication letter, explicitly suggesting an ``expanded version.'' Following this advice, I have substantially expanded the physical and statistical framework to provide the depth expected of a full ApJ article.

\textbf{Administrative Note on Publication Support:}

As an independent researcher, I would like to note for the editorial office that a 100\% waiver from the AAS Publication Support Fund was previously approved for the initial ApJL version of this manuscript (AAS76108) on April 24, 2026. As this new submission is the expanded realization of that exact same research, I respectfully request guidance on transferring or re-applying this approved support to the current ApJ submission.

I confirm that this manuscript is original, has not been published elsewhere, and is not currently under consideration by another journal.

Thank you very much for your time and assistance.

\closing{Sincerely,}

\end{letter}
\end{document}
